% Candidate only. The cited source statements and live-book interfaces are
% listed explicitly in SOURCE_LEDGER.md.
\begin{proof}
By the invariant-function correspondence (5.5) of \cite{BGJKK25}, there is
a unique finite continuous function $\varphi$ on $O$ such that
\[
                         h=\varphi\circ\operatorname{trop}.
\]
The test-form calculation in the proof of Proposition~5.3.10(iv) of
\cite{BGJKK25} applies to this continuous invariant function.  It identifies
$[\varphi]$ with $\operatorname{trop}_*[h]$, and the pushforward commutes
with $\dprimeddprime$.  The hypothesis therefore gives
\[
                         \dprimeddprime[\varphi]=0\quad\hbox{on }O.
\tag{*}\label{eq:p062-invariant-toric-boundary-zero-hessian}
\]

We show that $\varphi$ is locally smooth with zero Hessian on the tropical
toric partial compactification.  Fix $x\in O$, let $\sigma\in\Sigma$ be the
unique cone with $x\in N(\sigma)$, and choose a basic neighbourhood
$W=U(\sigma,\Omega,x_0)\Subset O$, where $\Omega\subset N_{\mathbb R}$ is
open, bounded, and convex.  As in the proof of
\cite[Theorem~3.1.4]{BGJKK25}, its dense part is
\[
                         W\cap N_{\mathbb R}=\Omega+\sigma.
\]
Restricting (*) to this open convex set shows that every distributional
second partial derivative of $\varphi$ is zero.  Thus
\[
                         \varphi(y)=\ell(y)+c
                         \qquad(y\in\Omega+\sigma)
\]
for a linear functional $\ell$ and a constant $c$.

For $y\in W\cap N_{\mathbb R}$ and $w\in\sigma$, the compactified ray
$[y,y+\infty w]$ lies in $W$.  Since $\varphi$ is finite and continuous at
its endpoint, $\ell(y+tw)+c$ has a finite limit as $t\to\infty$.  Hence
$\ell(w)=0$ for every $w\in\sigma$.  The affine function therefore descends
through $N_{\mathbb R}\to N(\sigma)$, and continuity plus density show that
it equals $\varphi$ on all of $W$.  Equivalently, $\varphi$ is constant
towards every boundary face of $W$.  Equation (3.11) of \cite{BGJKK25} then
says exactly that $\varphi$ is smooth for the standard Lagerberg structure,
and its smooth $\dprimeddprime$ is zero.

By Proposition~2.3.4 of \cite{BGJKK25}, the pullback
$h=\varphi\circ\operatorname{trop}$ is CLD-smooth on
$\operatorname{trop}^{-1}(W)$ and its smooth Hessian is zero.  We now make
the required PL and classical conclusion directly, without invoking any
additional smooth-zero-Hessian principle.  The lattice
$M\cap\sigma^\perp$ spans the dual of $N_{\mathbb R}/\langle\sigma\rangle$.
Choose $m_1,\ldots,m_r\in M\cap\sigma^\perp$ and real numbers
$\lambda_1,\ldots,\lambda_r$ such that
\[
                         \ell=\sum_{i=1}^r\lambda_i m_i.
\]
On the affine toric chart $X_\sigma$, which contains
$\operatorname{trop}^{-1}(W)$, both $\chi^{m_i}$ and $\chi^{-m_i}$ are
regular: $m_i,-m_i\in\sigma^\vee\cap M$.  Thus each $\chi^{m_i}$ is an
analytic unit there, and the defining convention for tropicalization gives
\[
 h=c+\sum_{i=1}^r\lambda_i\bigl(-\log|\chi^{m_i}|\bigr)
 \qquad\text{on }\operatorname{trop}^{-1}(W).
\tag{**}\label{eq:p062-invariant-toric-character-formula}
\]
Each summand in (**) is integral PL because its tropical representative is
the integral-affine function $m_i$; hence $h$ is locally $\mathbb R$-PL.

For every analytic unit $f$, both $\log|f|$ and $-\log|f|=\log|f^{-1}|$
are classically psh by \cite[Proposition~4.2(2)]{GR25Classical}.  The cone
property in \cite[Proposition~4.2(3)]{GR25Classical}, applied term by term
with the sign of $\lambda_i$, therefore makes both sides of (**) and its
negative classically psh.  Thus $h$ is classically pluriharmonic.  Since
$h$ is locally $\mathbb R$-PL on the good strict space $V$, the equivalence
in \cite[Theorem~4.3]{GR25Classical} makes both $h$ and $-h$ semipositive.
By Definition~7.2 of \cite{GJR26}, this is exactly Gubler--Jell--Rabinoff
harmonicity.  Since $x$ was arbitrary, all the stated conclusions hold on
$V$.
\end{proof}
